% =====================================================================
%  Georg Brandes, HANS BRØCHNER (1886).
%  In: Georg Brandes, Essays. Danske Personligheder (Kjøbenhavn: Gyldendalske
%  Boghandels Forlag (F. Hegel & Søn), 1889), pp. 230--255.  First printed in
%  Norden. Illustreret skandinavisk Revue (1887); reprinted in Samlede Skrifter II
%  (1899).  This is the 1889 text.
%
%  TRANSCRIPTION -- GENERATED; DO NOT EDIT BY HAND.  Rebuild from the repo root:
%    python3 ebook-method/hans-brochner/draft.py --emit --force
%    python3 ebook-method/collate/check.py hans-brochner --apply
%    python3 ebook-method/collate/check.py hans-brochner --mark --write
%    python3 ebook-method/hans-brochner/draft.py --post
%  Scan: ~/bibliotek/Brandes, Georg/danske-personligheder.pdf (SCANS.tsv,
%  brandes/hans-brochner), the Royal Library's scan, 250 ppi colour, ABBYY
%  text layer.  printed p = PDF page - 15 (p. 230 = PDF 245, p. 255 = PDF 270).
%
%  METHOD AND ACCURACY (2026-10-27).  E-book method (TRANSCRIPTION-PLAYBOOK §00, point 0); no one typed
%  the text.  The draft is Project Runeberg's text of this edition (runeberg.org/bgessays, pp. 230-255,
%  every page marked "proofread at least once"; Runeberg's page images appear to be this same KB scan),
%  built by ebook-method/hans-brochner/draft.py.  It was collated word by word with punctuation against
%  the KB scan's ABBYY text layer (ebook-method/collate/check.py; 4,887 words against 4,891): 19
%  disagreements, 1 settled by rule (a word division), 18 read at the image in BLIND crops (the reader
%  saw neither reading; two planted controls, both read right).  1 correction to Runeberg came of it:
%  „Jouffroy’s“ for Runeberg's „Jouftroy’s“ (p. 234).  Every other disagreement was the layer's fault.
%  Emphasis is italic (no letterspacing found: the layer has no runs of single letters, and the spot
%  pages and crops show none).  Every italic span Runeberg marks (10) and the 12 places where the stroke-
%  slant measure disagreed with it were read at the image: Runeberg's spans all stand, the measure was
%  wrong in every case (this face's italic leans little).  Two pages read whole at the image (pp. 232,
%  235, chosen by seed 20261027) agree with the text word for word, in paragraphs and in style.  An
%  unattested-word sweep (words found in no other transcription in the collection, 100) found no fault.
%  Page markers set at the exact turns by check.py --mark (26 of 26).
%  Word-accuracy: the whole text diffed against the KB ABBYY layer (check.py), 19 candidates,
%  1 transcription error corrected, 18 witness's fault, 0 unresolved.
%  What this cannot see: a misprint that Runeberg's proofreader and the ABBYY layer both silently
%  normalised (none was found on the two pages read whole).
%  Conventions: italics \textit; quotation marks «…» as printed; the apostrophe in „saa’“ (pp. 246, 247)
%  as printed; the verse quotation on p. 232 in small type, indented, as printed.  The head is printed as
%  title, date and a short rule; the essay ends with a double rule.  Running heads („H. Brøchner.“ /
%  „Essays.“) and sheet signatures are not transcribed.  The volume's Indhold lists the essay as
%  „H. Brøchner 230“.
%  PRINTED AS IS (found while translating, read at the image): p. 247 „fortabt i | i Glæden“ (i set twice
%  across a line end); p. 252 „vilde have frem;“ (a semicolon for a comma, in a quoted letter).
% =====================================================================
\documentclass[12pt,a4paper,openany]{book}
\usepackage[utf8]{inputenc}
\usepackage[T1]{fontenc}
\usepackage{libertinus}
\usepackage{libertinust1math}
\usepackage{textalpha}
\usepackage{graphicx}
\DeclareUnicodeCharacter{0254}{\reflectbox{c}}  % ɔ (U+0254) = reversed c
\usepackage[danish]{babel}
\usepackage[protrusion=true,expansion=true]{microtype}
\usepackage{setspace}
\usepackage[margin=1.3in]{geometry}
\usepackage{fancyhdr}
\usepackage{hyperref}
\hypersetup{pdftitle={Hans Brøchner (1886)}, pdfauthor={Georg Brandes}, hidelinks}
\setlength{\headheight}{15pt}
\pagestyle{fancy}
\fancyhf{}
\fancyhead[LE,RO]{\thepage}
\fancyhead[C]{\textit{Georg Brandes: Hans Brøchner (1886)}}
\setstretch{1.3}
\usepackage{marginnote}
\newcommand{\opage}[1]{\marginnote{\footnotesize\textit{[#1]}}}
% the essay's head as printed: title, date, a short rule
\newcommand{\essaytitle}[2]{\begin{center}{\large #1}\\[4pt]#2\\[4pt]\rule{1.5cm}{0.4pt}\end{center}}
% section numerals, centred
\newcommand{\sectionnum}[1]{\begin{center}#1\end{center}}
% the double rule printed at the end of the essay
\newcommand{\closingrule}{\begin{center}\rule{3cm}{0.4pt}\\[-0.8\baselineskip]\rule{3cm}{0.4pt}\end{center}}

\begin{document}
\thispagestyle{empty}
\begin{center}
\vspace*{2cm}
GEORG BRANDES:\\[2ex]
{\LARGE ESSAYS.}\\[2ex]
DANSKE PERSONLIGHEDER.\\[\fill]
\end{center}
\vfill
\begin{center}
KJØBENHAVN.\\[1ex]
GYLDENDALSKE BOGHANDELS FORLAG (F. HEGEL \& SØN).\\[1ex]
{\small GRÆBES BOGTRYKKERI.}\\[1ex]
1889.
\end{center}
\clearpage
\noindent\textit{Indhold} (p.~[v]), the entry for this essay: H. Brøchner \dotfill\ 230

\vspace{2em}
% ===== TEXT (pp. 230--255) =====
% --- p. 230 ---
\opage{230}\essaytitle{HANS BRØCHNER.}{(1886.)}

\sectionnum{I.}

Hvad man, nitten Aar gammel, med de
første Studie-Aar bag sig, søger hos Lærerne
ved et Universitet, det er ikke saa meget
Kundskab som Indtrykket af en kraftig, moden
Personlighed. Man attraar mindre Undervisning
end den Vejledning, som kan gives Ungdommen
af en ikke blot i Bøgernes Verden vel orienteret
Aand. Man ønsker sig en Lærer, der tillige
er en Karakter.

Der var i 1861 adskillige Yngre, som
glødede for Filosofi; næsten Alle troede paa
den. R. Nielsen søgtes af ikke Faa. Han
viste altid de Spørgende Interesse og Velvilje.
Men Brøchner tilbød udtrykkeligt i
Lektionskatalogen Vejledning for Enhver, der vilde
studere Filosofi. Det gjorde Udslaget for
En%
% --- p. 231 ---
\opage{231}kelte, hvem det var Alvor med en
sammenhængende Forskning. Og snart var det ikke
Kundskabstørst alene, der drog os til Brøchner.
Hans Personlighed drog.

Han var 41 Aar gammel, da vi saa ham
første Gang, hørte en Række Forelæsninger af
ham, fængsledes af hans elegante og strænge
Skikkelse, hans sydlandske Hoveds klassiske
Profil, det store alvorsfulde Blik under de
hvælvede, lidt trætte Laag. Han var, skjønt noget
medtaget af Aarene, den adeligste Mands-Skjønhed,
vi nogensinde havde set. Han var
høj og slank, mager og tør. Ansigtet var
smalt og noget indfaldent, det rige og bløde
Haar blaasort i sin Glathed. De store brune
Øjne havde et Udtryk, som havde de grædt
meget, som var deres Blik blevet blidnet af
Sorg. Panden hævede sig over dem stejl og
hvid, som i to Etager. Det spanske Skjæg
bidrog sit til at give Indtrykket af noget
Fremmed. Han lignede ikke en moderne Lærd,
snarere en Tænker fra Renaissancens Italien,
var da ogsaa efter eget Sigende af italiensk
Afstamning.

Han saa ud som en Mand, for hvem Vellevnet
ikke existerede; han satte øjensynligt sin
Ære i, ægte spinozistisk, ikke at have attraaet
eller at attraa jordisk Velfærd. Han levede i
% --- p. 232 ---
\opage{232}en filosofisk Fattigdom; thi selv da han var
gift, var hans Docentgage kun 2400 Kroner.
Hans Klædedragt bar Præget af streng Tarvelighed,
den var simpel og mørk, aldrig en Plet.
Vesten var som oftest imod Tidens Mode lukket
helt op til Halsen. Han bar Dag ud Dag ind
de samme Klæder, og de saa ens ud; han
syntes Intet at slide paa dem. Hans Ydre
gjenkaldte Welhavens Vers om den polske
Emigrant:

\begin{verse}\small
Hans Ansigt var falmet, hans Dragt var grov, \\
Hans Linned var rent og hans Hænder var smukke.
\end{verse}

Det var gjennem Øjet, ikke gjennem Øret,
at man fra først af følte sig tiltrukket. Thi
Foredraget var ængsteligt korrekt, pillent
soigneret i sin Sprogform, abstrakt i sin
Tankegang og besvarede intet Spørgsmaal anderledes
end gjennem en udførlig Fremstilling af
Spørgsmaalets og Svarenes Historie. Det bevægede
sig frem enstonigt, i afmaalte Sætninger, og
havde næsten ingen mundtlig Karakter, skjønt
det holdtes frit nok. Brøchner havde ikke
Talens Gave, som er den at bibringe hver
enkelt Tilhører det Indtryk, at Talen henvender
sig netop til ham. Hans Ord gik hyppigt bort
over vore Hoveder; de belærte os uden at
ramme os. Foredraget var følgerigtigt og
tanke%
% --- p. 233 ---
\opage{233}fast, men det lød, som læstes det op af en
indre Bog; det \textit{levede} ikke.

Taleren derimod levede for os. Det var
en Personlighed, som stod der, lige fornem og
interessant.

Han indgød Deltagelse; thi han saa ud,
som havde han engang været lutter Lidenskab,
medens han nu var lutter Tanke. Blikkets
sorgfulde Alvor tiltrak, paa samme Tid som
det noget Forlegne, Generte, der i Reglen laa
over Talerens Holdning, holdt i Afstand og
fremkaldte det Spørgsmaal, om man kunde
komme denne Mand menneskeligt nær.

Og han indgød Respekt. Det Lidt, vi
Unge vidste om hans Skjæbne, røbede en
ualmindelig Selvstændighed og Karakterstyrke. Han
havde, da han som ganske ung indstillede sig
til theologisk Examen, vist den Aabenhed og
det Mod at bekjende sin Opfattelse af
Kristendommens Dogmelære som mythisk. Som Følge
deraf var det blevet ham forbudt at tage denne
Examen. Han havde da kastet sig over Studiet
af Filosofi, havde oversat Strauss, havde i en
mod Martensen rettet Afhandling om \textit{Daaben}
støttet sig til Feuerbach, og var i den Grad
lagt for Had som farligt Menneske, at Christian
den Ottende egenhændigt havde strøget hans
Navn af Listen over Ansøgerne til et
Rejse%
% --- p. 234 ---
\opage{234}stipendium. Ikke desmindre havde han (ved
Bistand af en Broder) levet meget udenlands,
opholdt sig længe i Tyskland, været to Gange
i Italien og havde først efter 1848 paany
betragtet Danmark som et Land, der havde Plads
for ham.

Den Modsætning, hvori hans offentlige
Optræden stod til Rasmus Nielsens saa meget
livfuldere, men saa langt mere virtuosagtige
Foredrag, var slaaende for Alle. Det var en
Modsætning, der maa have svaret temmelig nøje til
den, som i Trediverne i Sorbonnen i Paris
Jouffroy’s Forelæsninger frembød til Cousins, hvis
ellers efterfølgende franske Skildring af Filosofen
Jouffroy, hvilken vi oversætte med nogle
Udeladelser, men uden at tilføje nogetsomhelst, har
stemmet med Virkeligheden:

«Man traadte ind i Jouffroy’s Forelæsningssal
og fandt der en mager Mand, lidt foroverbøjet
som alle Brystsyge, med dybtliggende Øjne,
med et tænksomt og melankolsk Fysiognomi og
med et Udtryk over hele sin Person af Træthed,
Adel og Resignation. Han var stolt,
forbeholden, gjerne taus, hverken imødekommende
eller fortrolig mod sine Elever. Medens Cousin,
naar han blev sat i Svingning af Begejstringens
Storm og det animalske Livs overstrømmende
Rigelighed, talte, meddelte sig, udøste sig,
for%
% --- p. 235 ---
\opage{235}klarede og overtalte med Haandbevægelser og
det hele oratoriske Tilbehør, ligegyldigt om i en
offentlig Have eller i sit Værelse, lige meget
hvor og lige meget for hvem --- var Jouffroy
selv paa Kathedret kold og behersket. Han
saa ud, som om han ikke vidste, at man
var der.

Hans Haandbevægelser var sparsomme,
hans Legeme ubevægeligt; man skulde tro, han
læste en usynlig Bog, ene og alene sysselsat
med at forstaa den, og kun vilde overbevise sig
selv; han tænkte højt. Ingen glimrende Ord,
ingen vovede Sætninger, ingen Bestræbelse efter
at more, vække eller røre; tvertimod lange
Indledninger, fulde af smaalige Inddelinger og
Under-Inddelinger, en omstændelig og uophørlig
Undersøgelse af de foreløbige Spørgsmaal.

Naar han endelig var trængt ind i sit
Æmne, ingen afgjørende eller forhastet
Sætning, uendelige Gjentagelser. Naar hans
Disciple gjennemlæste deres Optegnelser, fandt
de, at de havde nedskrevet den samme Idé
tre eller fire Gange. Og dog bøjede fra den
første Dag af enhver opmærksom Aand sig
under hans Aand. Man opdagede strax i ham
en skjult Arne, hvor der brændte en uudslukkelig
Ild; mere voldsom og mere mægtig end Cousins
blændende Illumination. Jo mindre han gav sig
% --- p. 236 ---
\opage{236}hen, desto mere følte man, hvor han brændte.
Disse fordoblede Gjentagelser angav ved deres
haardnakkede Forsøg en Aand, der var umættelig
efter det nøjagtige Udtryk, ude af Stand til at
slaa sig til Ro ved blot skitserede Tanker,
uovervindelig overfor Anstrengelser, udholdende
i sine Tankers Gang, indtil den havde naaet
fuldstændig Klarhed, endelig vel kjendt med
Vejens Sneverhed og med de Strømme af
Vildfarelser, der flyde paa begge Sider af den.
Man følte Ærbødighed og Tillid, og naar en
Skjælven i Stemmen eller en enkelt Gang et
Billede angav en vigtig Sandhed, han havde
fundet, iagttog man i dette svage Tegn mere
Bevægelse og mere Veltalenhed end i hans
Rivals glimrende Dithyramber.»

Det var neppe muligt at give en mere
træffende Skildring af Fortrinene ved Brøchners
Undervisning, om man udtrykkeligt gik ud
derpaa.

Og kritisk, som han fremfor Alt var anlagt,
glimrede han især som Polemiker. Hans
Angrebskrig mod Theologerne eller mod Nielsens
maskeret-theologiske Filosofi var et Mønster paa
overlegen Taktik. Han udviklede gjerne først
med største Nøjagtighed den Anskuelse, han vilde
bestride, fremsatte dernæst Punkt for Punkt
Beviserne for dens Uholdbarhed, anførte saa Alt,
% --- p. 237 ---
\opage{237}hvad der fra Modstanderens Side kunde tænkes
gjort gjældende mod disse Beviser, og
forklarede, hvorfor dette formodede Forsvar Intet
vilde betyde; saa beviste han, at andet Forsvar
end det anførte overhovedet ikke kunde føres.

Naar dette var gjort, opløste han Modstanderens
Begrebssammenhæng, gjorde ham Indrømmelser,
satte \textit{per impossibile} at et af de
kuldkastede Hovedpunkter var et fast Kastel, og viste,
at saa ligefuldt de mest afgjørende Kategorier
indeholdt Modsigelse paa Modsigelse. Skridt for
Skridt gik han rundt om sin Modstanders
Lærebygning, og vel bekjendt med dennes
Kampmaade stoppede han ethvert Hul, hvorigjennem
han kunde ville undslippe, indtil han havde
ham sikkert fangen. Saa endte han gjerne med
at rekapitulere alle de paaviste Modsigelser, og
sluttede tilsidst den i den alvorligste Tone holdte
Kritik med en Ironi, som man i al Sandhed og
uden Misbrug af Ordet kunde kalde frygtelig.

\sectionnum{II.}

Der gik blandt de unge Studenter en Legende
om, at Brøchner i sine unge Dage nede
i Italien havde været forlovet med en meget
fornem ung Dame, en Prinsesse eller Comtesse,
men Forbindelsen skulde være strandet paa
% --- p. 238 ---
\opage{238}Religionsforskjellen. Det gav ham en romantisk
Nimbus selv nu, da han var forlovet med en
ung dansk Pige fra Provinserne af borgerlig
Herkomst.

I 1862 blev han gift og levede i et lykkeligt
Ægteskab indtil 1866, da hans unge Hustru
døde efter at have født sit andet Barn. Der
var det bedste Forhold mellem de to
Ægtefæller trods en Aldersforskjel af en fuld Snes
Aar og tiltrods for det for en ung Kvinde
utilgængelige Omraade af abstrakt Filosofi, om
hvilket Brøchners Tanker og i Regelen ogsaa
hans Tale drejede sig. Hans Ægteskab var en
idyllisk men kort Episode.

I en meget længere Periode af hans Liv
stod Mindet om hin Italienerinde, Costanza T.,
hvem han i sin Ungdom havde kjendt og elsket,
som en Stjerne over hans daglige Existens.
Hun var bleven Drømmen i hans Liv, efter at
de to i Ungdommen vare blevne skilte ad. Et
meget smukt Daguerreotypibillede af hende som
ung hang altid i hans Arbejdsstue. Han talte
ikke sjældent om Originalen som om en ung
Dame, han engang havde kjendt, dog uden at
gaa nøjere ind paa deres Forhold.

I 1866 lærte jeg hende personligt at kjende.
Hun levede dengang, gift med en Landsmand af
sig, Grev Oreste B., i Paris, men i Grunden
uden%
% --- p. 239 ---
\opage{239}for Pariserlivet, i udelukkende Omgang med den
italienske, spanske og sydamerikanske Koloni.

Hendes Skjønhed, men især hendes tilsyneladende
Ungdom overraskede mig. Jeg havde
ventet at se en ældre Dame, men hun syntes
ikke mere end en sex eller syv og tyve Aar,
og hendes Skikkelse var ligesaa ungdommelig
som hendes Ansigt.

Man kunde neppe have nogen bedre
Anbefaling til hende end den fra Brøchner. Hun
kunde ikke høre nok om hans Liv og Færd.
Da hun i en første Samtale spurgte mig, om
Brøchner var bleven gammel, og jeg svarte, at
hans Haar endnu var sort, tilføjede hun: Og
hvilken sort Farve det har! Som aandelig
Individualitet existerede Brøchner egentlig ikke for
hende; hun var slet ikke trængt ind i hans
Væsen, havde end intet Forsøg kunnet gjøre
derpaa. Hun var et ungt Barn, da han lærte
hende at kjende, og var selv nu som udviklet
og meget fintfølende Kvinde i Grunden barnlig.
Hun var uvidende, som Italienerinderne er det,
yderst naturlig og ligefrem, men ikke desmindre
reserveret, altid i mildt og godt Humør, paa
Baggrund af en vis Vemod. Ikke just lykkelig,
men tilfreds. Hun var rolig, og hun syntes ved
Modsætningen til sin muntre, sprudlende,
verdslige og mere glimrende, men mindre smukke
% --- p. 240 ---
\opage{240}Søster, i hvis Hus hun opholdt sig, endnu mere
rolig og stilfærdig. I flere Maaneder saa jeg
hende dengang næsten dagligt; i 1870 har jeg
boet i hendes Hus i Landsbyen Alpignano ved
Turin, i 1871 atter truffet hende i Rom. Hun
havde den dybe Elskværdighed, der aldrig
formørkes af et Lune; der var den samme blide
Harmoni i hendes Sjæl, som laa over hendes
Træk.

Costanza var en forfinet Frugt af sit
Fædrelands yppige Jord. Hun lignede den ædle
Druesort, som man i Italien kalder Jordbærdruen,
der forener Druesaftens rige Sødme med
Skovjordbærrets naive Aroma. Man havde
handlet ilde med hende i hendes Stilling til
Brøchner. Hendes Moder, der var lige
forfængelig og bigot, havde efter Brøchners
Afrejse gjennem stadig Paavirkning i Hjemmet og
gjennem Skriftefaderens Indflydelse faaet
Forbindelsen hævet med den fremmede Kjætter.
Mange Aar efter blev hun saa gift med en
Fritænker, saa heftig som kun Italien fostrer dem,
langt voldsommere i sine Udtalelser end
nogensinde Brøchner, og knyttede sig af Hjertet til
ham, uden dog at istemme i hans Overbevisning.
Hendes Mand var en af de forstandigste,
retsindigste og dygtigst anlagte Italienere, jeg
% --- p. 241 ---
\opage{241}har kjendt, en Hader og Foragter af alt Skin.
Han vurderede Costanza, som hun fortjente.

Hun plejede stadigt Tanken om et Gjensyn
med Brøchner, og talrige vare de Planer, man
lagde til Rejser for ham i dette Øjemed. Men
Gjensynets Tanke syntes ikke at friste ham.
Han gik aldrig ind paa den, syntes snarere at
sky den. Det var neppe alene den Vanskelighed,
det for ham vilde være at foretage større
Rejser, som afholdt ham. Hverken Mangelen
paa Raad dertil eller hans svage Helbredstilstand
syntes det afgjørende. Rimeligvis skyede han
Sindsbevægelsen som smertelig. Eller han havde
vænnet sig saaledes til et Liv i Ungdomsmindet
og i Drømmeforholdet til Costanza, at han
frygtede Virkelighedens forstyrrende Pust. Han
maatte have Billedet af hende i Fjærnhedens
Afstand.

For hende som Kvinde og Sydlænderinde
existerede ikke denne Romantik. Hun havde
engang elsket ham, hun holdt inderligt af ham
endnu. Det vilde have været hende en Glæde
at gjense ham, at pleje ham, nu da han var
saa haardt syg. Hun var meget angst for, at
han nogensinde skulde fatte en ugunstig eller
urigtig Mening om hende. Da han giftede sig,
afbrød hun Brevvexlingen med ham, af Frygt
for at hans Hustru skulde finde noget Underligt
% --- p. 242 ---
\opage{242}i den. Men da saa Nytaar kom, greb hun
denne Lejlighed til at sende ham en
Lykønskning. Senere blev Brevvexlingen, dog med
store Mellemrum, optagen paany, og Costanza
tilbragte engang en Tid i ikke ringe Uro, da
hun under en haard og langvarig Sygdom havde
maattet lade sin Mand skrive for sig, og der
saa ikke fulgte noget Svar fra Brøchner. Hun
sendte ham derefter Evighedsblomster og røde
Negliker, hvilke sidste, efter en Ytring af
Brøchner, i det italienske Blomstersprog har «en rig
Betydning». I sine sidste Leveaar fortalte han
ikke sjældent til sine Smaabørn Anekdoter om
\textit{Costanza}, hvis Billede de kjendte. Det var
neppe uden en Tanke paa dette hendes Navn,
at han i Vilhelm Møllers «Nyt dansk
Maanedsskrift» undertegnede en Afhandling med
Pseudonymet Constans.

\sectionnum{III.}

Altfor lidet af Brøchners Væsen som
Menneske finder den almindelige Læser i hans
Skrifter. Indholdet er for abstrakt, og Stilen
skræmmer.

Paa et hængende Haar var han overhovedet
slet ikke blevet Forfatter. Han følte intet Kald
dertil; han havde modstandsløst ladet sin
Selv%
% --- p. 243 ---
\opage{243}kritik fortære sin Drift til Produktion, og var i
Fyrretyveaarsalderen naaet til det Standpunkt af
Ulyst ved al Frembringelse, at han glædede sig
eller troede at glæde sig, naar der andensteds, i
Tyskland f. Ex., blev fremsat Tanker, han
allerede for længe siden havde tænkt og
nedskrevet til eget Brug, men ikke givet i Trykken.
Hvad gjør det, sagde han i saadanne Tilfælde,
om det bliver sagt af mig eller en Anden, naar
det blot kommer frem.

23 Aar gammel havde han udgivet et polemisk
Lejlighedsskrift om Daaben; først fulde
14 Aar senere skrev han for at legitimere sig
til Ansættelse som filosofisk Docent ved
Universitetet, atter som Lejlighedsskrift, sin lille
Bog om Spinoza. Uden en yngre Vens ihærdige
Overhæng var det blevet derved. Men nu
fulgte, fra først af foranlediget ved den af Andre
begyndte Polemik mod R. Nielsen, efter 11 Aars
Taushed med korte Mellemrum fem store Bind
og adskillige mindre Arbejder, en hel Produktion,
nedskrevet efter Forfatterens 48de Aar, i
et Tidsrum, da han til det Yderste var svækket
af den Lunge-Svindsot, der ti Aar igjennem
nedbrød ham Tomme for Tomme, knækkede
ham, som ved den skærpede Dødsstraf, Led
for Led.

Men selve denne Sygdom, der ikke gav
% --- p. 244 ---
\opage{244}Pardon, understøttede Opfordringerne udefra.
Saalænge Brøchner endnu var rask eller dog
ikke troede sig livsfarligt angrebet, længtes han
mangen Gang mod Døden. Saa virksom han
end var, saa havde han ikke alene ingen Frygt
for Døden, men en Art Kjærlighed til den.
Han kunde tale om den med Varme. Den
betød for ham Befrielse fra Personligheden,
Gjenforening med hint Allivets upersonlige Princip,
fra hvilket Individet er en Udstrømning; og han
var fuldt tilfreds med sin Selvforvandling til et
Moment i Guddommens Viden af sig selv i
dens evige Aandsrige --- som han mystisk
udtrykte det.

Da Døden nærmede sig med hurtige Skridt,
fik den vel ikke derved i og for sig nogen
afskrækkende Karakter for ham; men han stræbte
at leve saa længe som muligt, dels for at kunne
lægge Grundvolden til sine Børns Opdragelse i
den Aand, hvori han ønskede dem opdragne,
dels for gjennem Offentliggjørelse af sine
Manuskripter at efterlade Vidnesbyrd om, hvad han
havde tænkt og forsket. Saaledes blev hans
sidste Leveaar hans eneste literært produktive.

Men disse hans store Bøger afspejle kun en Del
af hans Menneskevæsen og ikke den menneskeligste.
De give en Forestilling om hans Tankeliv
og hans pantheistiske Religiøsitet. De røbe hans
% --- p. 245 ---
\opage{245}Skarpsindighed og hans Pathos --- det er alt.
Og, som jeg sagde det, Stilen skræmmer. Det
er en rent skolastisk Stil, et Bogsprog med en
Mangfoldighed af tyske Participialkonstruktioner,
hvor Alt er \textit{ordentligt}, men Intet har Livets
Uvilkaarlighed. Den afspejler alt det Stive, det
Nøjeregnende og ængsteligt Punktlige i
Brøchners Natur, giver en Slags Forestilling om en
vis Smaalighed i Smaating, der hos ham ikke
udelukkedes af hans Grundanskuelses Storsyn.
Men der findes en helt anden Stil i hans Breve
og i hans faa sidste populære Ting. Der er
over denne stive Stil i de store Bøger et Skjær
af Pedanteri, som aldrig laa over Brøchners
Person. Man ser ret vel, at han var en streng
Metafysiker, men man sporer ikke hans Stilling
til den konventionelle Moral, hans frie og vide
Tolerance, hvor noget af Verdensmandens
Synsmaade var forenet med noget af Spinozas høje
Ligegyldighed. Man formaar ikke at lodde
Dybden og at maale Styrken af hans Afsky for
al Theologi og alle theologiske Opfattelser, stillet
overfor denne hans Tankes officielle Kathederform,
der er saa videnskabelig i sit Tilsnit og
saa lidenskabsløs i sit Udtryk. Man ser ikke
heller, hvor pirrelig, hvor redebon til pludseligt
Brud han var overfor de Omgangsfæller, med
hvem ingen dybere Sympathi forbandt ham.
% --- p. 246 ---
\opage{246}Og hvem vilde endelig af disse Værker kunne
slutte sig til, hvad Brøchner var som Ven! Var
han øm og varigt hengiven som Elskende, mønsterværdig
som Fader, der maatte være Smaabørn
i Moders Sted, saa var han dog i endnu højere
Grad sjælden som Ven. Han, der var forbeholden
indtil Generthed, mundtlig altid spotsk
overfor ethvert Udbrud af ungdommelig
Enthusiasme, var lidenskabelig i Venskab. Overfor
meget yngre Venner var han paa én Gang som
en Ven og som en Fader. Han tænkte uafbrudt,
Aar ud, Aar ind, med en Omsorg, der
ikke forbigik nogen Detail, paa hvad han
formaaede at udrette for Vennen. Hans Venskab
var aktivt, en stadig Virksomhed. Og hvad
der er ejendommeligt, det tog fra Aar til Aar
bestandig til i Styrke. Det steg i Inderlighed
og forstaaende Kjærlighed og naaede --- maaske
efter fjorten-femten Aars Forløb --- sin største
Højde i det Aar, han døde.

Men dette hænger sammen med det Grundtræk,
som ogsaa kun vanskeligt fornemmes
gjennem de store Bøger, at Brøchner var i
stærk Udvikling, næsten lige til Opløsningen
kom.

Det saa’ ikke saadan ud. Han havde
længe været en ensom Mand. Han kjendte
altid kun faa, og han havde i de sidste Aar
% --- p. 247 ---
\opage{247}trukket sig stedse mere tilbage fra Verden, var
ikke heller af dem, der paavirkes synderligt af
deres Omgivelser. Han havde ogsaa forholdsvis
tidlig faaet noget af den ældre Mands Væsen.
Hans Livsindtryk laa fæstnede i en Sum af
Anekdoter, som han tidt vendte tilbage til uden
at huske paa, han havde fortalt dem før. Han
havde aldrig været nogen smidig Aand, og hans
Antipathier, de literære og videnskabelige som
de personlige, var af urokkelig Natur. Hans
Ringeagt for en Tænker som Schopenhauer
oversteg alt fornuftigt Maal; hans Afsky for
den Form af national Selvforgudelse, han fandt
hos Grundtvig, og for den Blanding af Dilettanteri
og Uvederhæftighed, han mente at finde
hos Goldschmidt, kjendte ingen Grænse.
Allerede naar man i Tredserne saa’ ham sysle i sin
Have i Søgade, optaget af at udvikle stedse
nye Varieteter af Stedmodersblomster, fortabt i % PRINTED AS IS ("i" set twice, at the end of one line and the start of the next)
i Glæden over de talrige forskjellige
Farvemønstre, kunde han gjøre Indtrykket af en Mand,
der var færdig med Menneskelivet og som fra
det havde tyet over til Samliv med Planten.

Men det var et Skin. Han undergik i de
senere Aar en kjendelig Metamorfose. Hans
tidlige Ungdoms Kamplyst var vaagnet i ham
paany. Hans Tilbageholdenhed var bleven
ringere, og hans politiske og sociale Anskuelser
% --- p. 248 ---
\opage{248}antog en ny Farvning. Han fjernede sig helt
fra Personligheder, hvem han tidligere i politiske
Sympathier havde staaet nær; mellem ham og
hans gamle Partifæller, de Nationalliberale, skete
i hans Sind det mest afgjørende Brud; han stod
Socialisterne langt nærmere end dem.

Han havde aldrig været tilbøjelig til
Overvurdering af Forholdene i Danmark; men i
disse sidste Aar, da hans Sjæl følte Driften til
at røre sine Vinger, sporede han stærkere
Lidenhedens og Smaalighedens Tryk. I et Brev fra
1870 siger han: «Hvad der her hjemme fra er
at fortælle Dem er saa lidt og saa fattigt, at
det neppe er Omtalen værd. Vort offentlige
Liv gaar sin trivielle Rigsdagsgang, med liden
Gjerning og megen Tale, helst om hvad der
ikke er Tale værd; med smaa Ideer og smaa
Personligheder, men med store Illusioner, stor
Forfængelighed og store Smaaligheder. Det er
ikke glædeligt, men det maa vel saa være som
Gjennemgang til et virkeligt frit Liv.»

I et Brev fra 1874 hedder det endnu
smerteligere og ligesom med et vemodigt
Tilbageblik: «Det er fysisk usundt at leve i et
lille beklumret Kammer; det er aandeligt usundt
at leve i saadanne Dværgforhold som herhjemme.
Alt skrumper saa ynkeligt ind; man faar aldrig
det friske Indtryk af det Store; man mærker
% --- p. 249 ---
\opage{249}ikke Verdenstankelivets kraftige Pust. Hvor let
synker man ikke sammen herhjemme! hvor let
tager man ikke til Maalestok for sin Stræben
det Smaalige, der omgiver En! hvor let standser
man ikke i Løbet, naar man mærker, at man
er et Stykke forud for de Faa, der kæmpede
med om Prisen!»

I disse Aar gik det endelig op for ham,
at han ved en anden Skrivemaade kunde vinde
sig en større Læsekreds. Han gik imellemstunder
endog saa vidt som til at beklage, at
han ikke som ung havde opholdt sig længe
nok i Frankrig til at modtage en Paavirkning
af fransk Fremstillingskunst --- en Ytring, der i
hans Mund maatte forbause, eftersom han hidtil
slet ikke selv havde havt nogen Følelse af
Manglerne ved hans Bøgers Form.

Han havde hidtil ligesom ikke vovet at
skrive med hele sit Menneskevæsen. Hans
Brochure mod Nielsen var det første Symptom
paa hans Drift til at anslaa en naturligere Tone.
Dog smelter i den det Vittigtskarpe ikke ret
sammen med det Højtidelige. I den lille
Afhandling om Giordano Bruno er der en tydelig
Bestræbelse efter at male, stærke Tilløb til en
helt ny Fortællemaade; og i de efterladte
\textit{Erindringer om S. Kierkegaard} er Fremstillingen
i al sin Fordringsløshed ypperligt skarpttegnende.
% --- p. 250 ---
\opage{250}Hvert Ord her gjengiver Virkelighed; selv det
Ubetydeligste har Liv.

Han behøvede kun at anvende den utvungne
Stil, han altid havde havt til sin Raadighed i
mundtlig Meddelelse og i Breve, men som han
ikke havde turdet fortro sig til.

Jeg havde rettet en Del Indvendinger imod
hans Hovedværk \textit{Bidrag til Opfattelsen af
Filosofiens historiske Udvikling.} Jeg havde paavist
Ensidigheden af en Fremstilling, der kun gav
Filosofiens Længdesammenhæng, og Inkonsekvensen,
som laa i, at Brøchner undertiden tog
Bredesammenhængen med, undertiden ikke. Naar
Talen er om den græske Filosofi, er han villig
til at se den i dens Sammenhæng med den hele
græske Aand og Kultur og mener, at den kun
saaledes kan forstaaes. Men naar Talen er om
moderne fransk, tysk, engelsk Filosofi, da
behandles den, som om den førte en isoleret
Existens uden Forhold til Folkeaanden, uden
Forhold til Landets Historie, Poesi o. s. v. Jeg
udtalte min Beundring for hans Gave til at gribe
Hovedpunkterne og fremstille Hovedtrækkene i
et System; men jeg spurgte: Er hos alle Filosofer
Systematiken det Væsentligste og Vigtigste?
Er Skelettet Alt, er Legemet Intet? Bør man
ikke idetmindste frygte for at stemple
interessante uensartede filosofiske Fænomener med
% --- p. 251 ---
\opage{251}samme officielle Skoleord, naar rent forskjellige
Nationalaander er komne tilorde gjennem dem?
Condillac fremstilles gjerne som et Appendix
til Locke; er der ikke ligefuldt et gabende
Svælg mellem deres hele Maade at filosofere
paa? Er der ikke paa Nationalitetens Vegne
et endnu stærkere Slægtskab mellem Descartes
og Condillac, end der med Hensyn til Tænkningens
Resultater er imellem Condillac og Locke?

Kort sagt, jeg skjulte ikke, at den
Behandlingsmaade, der stod for mig som den
fyldestgjørende af Filosofiens Historie, var en ganske
anderledes konkret og levende end Brøchners.
Han svarede med en Henvisning til den
Langsomhed, med hvilken Nationaliteten hos de
Nyere har arbejdet sig frem af Latinens
Fællesskab, og til de saa yderst sammensatte Forhold
hos de Moderne. Men havde blot hans Bog
været skrevet i den Stil og med det Liv, der
lægger sig for Dagen i det følgende Brevsted:

«Den græske Filosofi var ægte national og
kunde derfor ikke skyde \textit{frugtbare} Skud i en
fremmed Jordbund. Den nyere udfolder sig paa
et større Terræn; den kaster som det indiske
Figentræ Grene ud til alle Sider, der bøje sig
mod Jorden, slaa nye Rødder, og voxe med
frodig Væxt. Derfor bliver Complicationen saa
stærk, Paavisningen af Udviklingen saa forviklet,
% --- p. 252 ---
\opage{252}at de store Træk, jeg vilde have frem; % PRINTED AS IS (a semicolon where the sense wants a comma)
let blive
utydelige. Prøv engang at analysere Spinoza fra
Nationalitetens Synspunkt. Prøv at bestemme
Blandingen af det Jødiske, det Europæisk-Sydlandske,
det Hollandske og det universelt Moderne.
Prøv at analysere Kant med hans Aandsslægtskab
med Tyskerne, hans Blodsslægtskab
med Skotterne, hans Valgslægtskab med
Franskmændene. De vil se, Opgaven ikke er let; jeg
er bange for, at den bliver lige saa vanskelig
som at beregne de indbyrdes Perturbationer af
tre Himmellegemer.»

Jeg mener ikke desmindre, at dette var
Opgaven, det gjaldt om at løse, og det
forekommer mig, at blot den Maade, paa hvilken
Brøchner her tumler med den, viser, at han ved
en noget forskjellig Udvikling af sine Kræfter
vilde have formaaet at lægge Sten til den
Bygning, han betegner som uopførlig.

\sectionnum{IV.}

En Ugestid før sin Dødsdag (han døde 17.
December 1875) tog Brøchner Afsked med sine
Venner. Da jeg en Dag traadte hen til hans
Sygeleje, gav han mig en tyk Pakke i Haanden ---
det var alle de Breve, jeg i en lang Aarrække
havde skrevet til ham, ordnede efter Tidsfølgen
% --- p. 253 ---
\opage{253}og paa Bagsiden af Konvoluten forsynede med
Afsendelsens Dato, Modtagelsesdatoen og den
Dato, da Brevet var besvaret. Hans stærke
Ordenssans og hans Omhu for at ingen fortrolig
Meddelelse fra en Ven skulde falde i fremmede
Hænder, røber sig i dette lille Træk.

Han havde en kort Samtale med Enhver,
der kom, og trykkede kort og kraftigt hans
Haand til det endelige Farvel.

Dette sidste Sygeleje, som afsluttede en
Lidelseshistorie, der strakte sig over mere end
et Decennium, var yderst kvalfuldt.

Alt som Tanken droges hen til den
nærforestaaende Død, kom den Syge til at sysle
med Forestillingen om sin Begravelse. Tanken
om Kirken, til hvilken han intet Forhold havde
havt, og om dens Præster, der altid havde
været ham imod, var ham meget pinlig. Og i
sin Følelse af helt at have tilhørt Videnskaben
og ingen Berøring at have havt med Gejstligheden,
formulerede han dette Ønske, som han
fik rette Vedkommende forebragt, at blive jordet
fra Universitetets Forsal og ikke fra det
nørrebroske Kapel. Han udtalte dernæst Ønsket om,
at en yngre Ven vilde tale over ham istedenfor
Præsten, og nævnte den, hvem han dertil havde
udset, og som ogsaa havde været villig dertil.

Naturligvis kom der efter hans Død et
af%
% --- p. 254 ---
\opage{254}slaaende Svar med Hensyn til Jordfæstelsen fra
Universitetets Aula, og enhver Tanke om den
paagjældende Vens Optræden ved denne Lejlighed
faldt dermed bort.

Skjøndt Brøchner havde den mest levende
Uvilje mod, at en Præst skulde tale over hans
Kiste --- han sagde paa sit Dødsleje: Jeg tror,
at hvis saadan en sort Mand begynder at præke
over mig, rejser jeg mig op af Kisten og
protesterer --- blev den gejstlige Tale ikke undgaaet.

Han blev begravet den 23. December.
Man havde, synes det, bibragt den Gejstlige
den Forestilling, at Brøchner havde ønsket hans
Medvirkning; han antydede det og udtalte bl. A.
sin Glæde over, at den Afdøde og han, skjøndt
adskilte ved deres Overbevisninger, vare mødtes
i gjensidig Agtelse. Næste Dag meddelte et
dengang meget læst Blad med Aplomb, at
Præsten havde talt efter Brøchners udtrykkelige
Ønske.

Iøvrigt blev han efter sin Død skildret med
netop den Velvilje og Varme, som udfordredes
til at vise, hvor stor Afstanden var mellem en
\textit{Idealist} som ham og andre yngre Fritænkere af
helt forskjellig Støbning. Ja selve dette, om
han med Rette kunde kaldes Fritænker, er
efterhaanden blevet et stedse større Spørgsmaal.

Et Spørgsmaal er det ikke. Men om
Aands%
% --- p. 255 ---
\opage{255}sneverhed vilde det vidne, hvis man ansaa
Brøchners Personlighed for udtømt i hint golde Ord.
Vore Meninger ere ikke den væsentligste eller
oprindeligste Del af vort Væsen. Brøchner var
af Naturel en Randsager, hvis lidenskabelige
Forskning sattes i Bevægelse og dreves af en
\textit{Foraarsvilje} som den, der driver sit Spil i
Naturen, naar Safterne stige i alle Planter. Han
var som Aand en Grubler og en Naturfrom,
hvis Tro paa Menneskeaanden og hvis Forhold
til Universets Aand havde Karakteren af en
Religion. Han staar i dansk Literaturs Historie
som en ny Tids Bebuder.

\closingrule


\end{document}
